\documentclass{article}
\textheight 23.5cm \textwidth 15.8cm
%\leftskip -1cm
\topmargin -1.5cm \oddsidemargin 0.3cm \evensidemargin -0.3cm
%\documentclass[final]{siamltex}

\usepackage{ctex}
\usepackage{verbatim}
\usepackage{fancyhdr}
\usepackage{graphicx}
\usepackage{booktabs}
\usepackage{amsmath}


\title{hwcode3实验报告}
\author{孙天阳 SA23001051}
\begin{document}
\date{2024 年 10 月 19 日}
\maketitle

利用有限元方法求解两点边值问题
\[
\begin{cases}
-u'' = f, \quad 0 < x < 1,\\
u'(0) = u'(1) = 0.
\end{cases}
\]
其对应的变分问题为寻找$ u\in H^1([0,1]) $满足
\[a(u,v)=\int_0^1u'v'\mathrm{d} x=(f,v)=\int_0^1fv\mathrm{d} x,\quad \forall v\in H^1([0,1]).\]
在本次作业中用来近似$ H^1([0,1]) $的是连续分段线性多项式空间.具体而言,在$ [0,1] $中等距插入$ N $个内部节点,将之划分为$ N+1 $个区间$ [x_{i-1},x_{i}] $, $ i=1,\cdots,N+1 $.在每个区间$ [x_{i-1},x_{i}] $上, $ u $是一个线性函数,因此$ u $在$ [x_{i-1},x_{i}] $上的表达式可以由两个端点的插值决定
\[u=f^i_0(x)u_{i-1}+f^i_{1}(x)u_i,\]
其中$ f^i_0(x),f^i_{1}(x) $是Lagrange插值多项式.定义形函数
\[\vec{N}^i(x)=\begin{pmatrix}f_0^i(x)&f_1^i(x)\end{pmatrix}\]
定义$ \vec{u}_i=\begin{pmatrix}u_{i-1}&u_{i}\end{pmatrix}^T $,则$ u $在$ [x_{i-1},x_i] $上就是$ \vec{N}^i(x)\cdot \vec{u}_i $.再算一下导数
\[u'(x)=\vec{B}^i(x)\cdot \vec{u}_i,\quad \vec{B}^i(x)=\frac{\mathrm{d}}{\mathrm{d}x}\vec{N}^i(x).\]
这样一来
\[a(u,v)=\sum_{i=1}^{N+1}\int_{x_{i-1}}^{x_i}u'v'\mathrm{d} x=\sum_{i=1}^{N+1}\int_{x_{i-1}}^{x_i}\vec{u}_i^T\vec{B}^i(x)^T\vec{B}^i(x)\vec{v}_i\mathrm{d} x=:\sum_{i=1}^{N+1}\vec{u}_i^TA^i\vec{v}_i=\vec{u}A\vec{v},\quad (f,v)=F\vec{v}\]
注意目前$ \vec{u} $是一个$ N+2 $维（2个端点和$ N $个内部节点）的向量,有线性方程组
\[A\vec{u}=F.\]
我们有一个观察,如果$ u $是原微分方程的解,那么$ u+c $也是原方程的解,即如果$ \vec{u} $是线性方程组的解,那么$ \vec{u}+\vec{c} $也该是线性方程组的解,即$ A\vec{c} $应该为$ 0 $.所以理论上$ A $不是一个满秩矩阵,我们得不到唯一解$ \vec{u} $,所以我们通过多加一个约束的方式消除掉额外的这个自由度使得它跟我们取的准确解对齐,这里选择对齐二者的积分平均值,使用梯形法近似积分,也就是要求
\[\frac{u_0}{2} + u_1 + u_2 + \cdots + u_N + \frac{u_{N+1}}{2} = \int_0^1x^2(x-1)^2\mathrm{d} x\]
因为$ A $奇异,在matlab中直接使用$ A\backslash F $求解线性方程组会报错,所以我们将矩阵$ A $的最后一行替换为
\[\begin{pmatrix}\frac{1}{2}&1&\cdots&1&\frac{1}{2}\end{pmatrix}.\]

\end{document}
